\section{Related Work}
\subsection{Cloud Removal Methods (Overview)}
Cloud removal methods in optical remote sensing are commonly grouped into three categories: (i) single-image methods, (ii) multimodal-based methods, and (iii) multitemporal-based methods~\cite{CR_Advances_Review_ORS}.

\textbf{Single-image methods} aim to reconstruct surface information using only the cloudy optical image and typically rely on statistical or physical priors (e.g., frequency filtering or atmospheric scattering models). Zhang et al.~\cite{single_variation} proposed the Haze Optimized Transformation (HOT) for thin cloud and haze compensation, while He et al.~\cite{single_haze_removal_dark_prior} introduced the dark channel prior, later adapted for thin cloud removal in optical remote sensing. With deep learning, CNNs, U-Nets, and GAN-based architectures have also been applied in this setting; however, single-image methods remain fundamentally limited under dense cloud cover due to missing surface information~\cite{CR_Advances_Review_ORS}.

\textbf{Multimodal-based methods} integrate auxiliary data from complementary sensors to improve restoration. In practice, the most significant progress has been achieved by fusing SAR and optical imagery, leveraging the cloud-penetrating capability of SAR. Meraner et al.~\cite{CR_SEN2_dRNN} proposed DSen2-CR, combining Sentinel-1 and Sentinel-2 to enhance reconstruction under thick clouds and preserve spectral fidelity. Grohnfeldt et al.~\cite{A_cGAN_fuse_sar_MS_CR} demonstrated the effectiveness of conditional GANs for SAR–optical fusion, and Xu et al.~\cite{GLF_CR} proposed GLF-CR using a global--local fusion strategy. This line of work motivates SAR-to-optical translation approaches, which either translate SAR features into optical-like imagery or fuse SAR with partially corrupted optical inputs, while still facing challenges related to data registration, modality differences, and SAR speckle.

\textbf{Multitemporal-based methods} exploit repeated observations to fill cloud-covered regions. Xu et al.~\cite{CR_spars_repre_MT_dict_L} proposed multitemporal dictionary learning for reconstructing thin and thick cloud regions without explicit cloud masks. Ebel et al.~\cite{UnCRtainTS} introduced UnCRtainTS, combining multispectral reconstruction with per-pixel uncertainty estimation. While highly effective under persistent cloud cover, multitemporal approaches can be sensitive to geometric misalignment and temporal land-cover changes~\cite{expl_ML_CR_Cameroon}, and they require consistent time series data; mono-temporal approaches avoid such requirements and reduce co-registration constraints~\cite{kiwa_auto_Delineation_BAs}. Given the focus of this work on mono-temporal SAR--optical generation, the following subsections review SAR-to-optical translation methods and their use for cloud removal.

\subsection{SAR-to-Optical Image Translation}
SAR-to-optical translation is defined as an image-to-image (I2I) task that synthesizes optical-like imagery from SAR data, despite the domain gap between backscatter intensity and spectral reflectance as well as SAR-specific artifacts such as speckle and geometric distortions. Before the emergence of generative AI, SAR-to-optical translation relied on heuristic or classical methods. Early approaches used pseudo-colorization of SAR channels or polarization composites to enhance interpretability~\cite{colorization_Schmitt_2018}, though they failed to reproduce true optical characteristics. Multisensor fusion methods, such as combining SAR with prior cloud-free optical data via intensity–hue–saturation (IHS) or wavelet-based transforms~\cite{IHS_wavelet_Zhang2019}, provided partial solutions but depended on handcrafted features and complex preprocessing, limiting scalability and accuracy.

Recent deep learning approaches—most notably generative adversarial networks (GANs)~\cite{GANs_Goodfellow} and diffusion models~\cite{DDPM_2020}—have enabled end-to-end learning of mappings from SAR to optical observations by modeling cross-modal feature relationships and pixel-level distributions. Most existing studies focus on reconstructing the visible (RGB) subset of Sentinel-2 (or extend to near-infrared), whereas comparatively fewer works address full multispectral reconstruction; correspondingly, the literature often distinguishes SAR-to-optical translation from SAR-to-multispectral (SAR-to-MS) translation. In this paper, the term ``optical'' denotes the full 13-band Sentinel-2 spectral range, and the proposed methodology extends SAR-to-optical translation to the complete multispectral domain.


The foundational GAN framework by Goodfellow et al.~\cite{GANs_Goodfellow}, later extended to Conditional GANs (cGANs) by Mirza and Osindero~\cite{cGANs_mirza}, enabled image-to-image translation tasks that made SAR-to-optical synthesis feasible. Fuentes Reyes et al.~\cite{sar2opt_cGAN_Optim_oppr_limits} optimized an unsupervised CycleGAN for SAR-to-optical translation, improving interpretability and reducing speckle through customized preprocessing and architectural refinements, although fine structural details remained difficult to preserve in urban scenes. Wang et al.~\cite{s2ot_s_cycle_adv} introduced a Supervised CycleGAN (S-CycleGAN) by incorporating a pixel-wise MSE loss, producing realistic optical images and showing strong potential for cloud removal. Gao et al.~\cite{cr_fuse_HR_GEN} extended this idea with a fusion-based GAN framework for high-resolution imagery. Instead of directly translating SAR to optical, they generated a simulated optical image from SAR data and fused it with both SAR and optical inputs to reconstruct more realistic results. An ablation study confirmed that this two-stage fusion strategy achieved superior performance among tested configurations.

Following the success of Vision Transformers (ViTs)~\cite{ViT_2020}, transformer-based architectures have also emerged. Zhao et al.~\cite{hvt_cgan} proposed the Hybrid Vision Transformer cGAN (HVT-cGAN), combining CNN and ViT branches to capture both local details and global semantics. A Convolutional Attention Fusion Module (CAFM) adaptively merged multiscale features, enhancing texture and color fidelity. Trained on the SEN1-2 dataset~\cite{sen12_2018}, HVT-cGAN achieved superior visual and quantitative results over previous GAN-based models. Park et al.~\cite{s2o_ViT_cGAN} further improved this approach with a multiscale ViT-based cGAN architecture integrating perceptual loss and a two-phase transfer learning strategy to enhance realism and stability.

Diffusion models have recently entered this field. Bai et al.~\cite{c_diffusion_s2o} introduced a conditional diffusion model for SAR-to-optical translation, which, despite limited experimentation due to computational constraints, demonstrated promising performance compared to GAN-based methods. Bai et al.~\cite{s2o_color_super_diff} later extended the framework with color supervision to improve reconstruction fidelity. A more recent contribution, the Multi-Temporal Conditional GAN (MTcGAN) by Kwak and Park~\cite{assessing_MT_cGANS_s2o_crop}, was designed for early-stage crop monitoring and utilized SAR–optical pairs from reference and prediction dates to capture temporal dynamics, achieving superior spectral consistency compared to conventional methods.

These developments have motivated the use of SAR-to-optical synthesis as a data-driven strategy for cloud removal, discussed next.

\subsection{SAR-to-Optical Image Translation for Cloud Removal}
SAR-guided cloud removal exploits the cloud-penetrating capability of SAR to reconstruct missing optical observations in cloud-contaminated scenes. Early methods relied on traditional signal processing techniques: Huang et al.~\cite{huang2015} introduced sparse representation-based cloud removal using SAR data, which Xu et al.~\cite{CR_spars_repre_MT_dict_L} extended via multi-temporal dictionary learning. These approaches demonstrated the feasibility of SAR-assisted cloud reconstruction but struggled under heavy cloud cover or rapidly changing surface conditions.

With deep learning, multimodal fusion approaches became dominant. Enomoto et al.~\cite{filmy_cloudy_NIR_Enomoto} applied a cGAN for cloud removal by fusing the RGB composite of a cloudy optical image with the cloud-free near-infrared (NIR) band to reconstruct a cloud-free RGB image. Although limited under dense cloud conditions, the approach was pivotal in demonstrating the potential of multimodal data fusion for cloud removal and laid the groundwork for subsequent studies. Building on this concept, Grohnfeldt et al.~\cite{A_cGAN_fuse_sar_MS_CR} introduced SAR-Opt-cGAN, a model designed to fuse Sentinel-1 SAR and Sentinel-2 optical data—marking the first use of SAR data within a cGAN framework for cloud mitigation. Their adaptation of the Pix2Pix architecture~\cite{pix2pix_2018} allowed flexible multi-channel inputs and was trained on a subset of the SEN1-2 dataset~\cite{sen12_2018}. The results confirmed the advantage of incorporating SAR data, validating the effectiveness of the approach for mitigating cloud cover.

Beyond direct cGAN fusion, specialized pipelines were proposed. Darbaghshahi et al.~\cite{CR_RS_GAN_s2o} introduced a dual-GAN framework with Dilated Residual Inception Blocks (DRIBs), translating SAR to optical in the first stage and fusing with cloudy optical inputs in the second; it showed strong qualitative results but moderate quantitative accuracy. Ebel et al.~\cite{UnCRtainTS} advanced this direction with UnCRtainTS, which reconstructs multispectral imagery while predicting per-pixel uncertainty, improving performance and providing reliability indicators.

Recently, diffusion models have become highly competitive. Zou et al.~\cite{DiffCR} proposed DiffCR, a fast conditional diffusion framework and reported state-of-the-art on SEN12-MS-CR~\cite{sen12ms-cr_2021}. DiffCR employs a decoupled conditional architecture and a Time and Condition Fusion Block (TCFBlock), and predicts cloud-free images directly, enabling high-quality results in as few as 1–5 denoising steps with substantially reduced computational cost. In parallel, Meraner et al.~\cite{CR_SEN2_dRNN} proposed DSen2-CR, a deep residual SAR–optical fusion network with a Cloud-Adaptive Regularized Loss (CARL) to preserve uncorrupted content while focusing reconstruction on clouded regions; evaluated on SEN12-MS-CR, it showed strong global generalization and high reconstruction accuracy across all 13 bands, with best performance on surface-related bands and higher errors on atmospheric bands.

Overall, the literature demonstrates a progression from classical SAR-guided reconstruction to deep multimodal fusion, uncertainty-aware prediction, and diffusion-based models, establishing SAR-to-optical translation as a strong paradigm for cloud removal. This work builds on that foundation by evaluating SAR-to-optical translation for cloud removal across the full 13-band Sentinel-2 spectrum and analyzing model behavior through bandwise evaluation and ablation studies.
